\documentclass[11pt,a4paper]{article}

% ---------- Packages ----------
\usepackage[utf8]{inputenc}
\usepackage[T1]{fontenc}
\usepackage{lmodern}
\usepackage[french]{babel}
\usepackage{amsmath,amssymb}
\usepackage[margin=2.3cm,headheight=15pt]{geometry}
\usepackage{fancyhdr}
\usepackage{enumitem}
\usepackage{xcolor}
\usepackage[most]{tcolorbox}
\usepackage{hyperref}
\usepackage{tikz}

% ---------- En-tête / pied de page ----------
\pagestyle{fancy}
\fancyhf{}
\fancyhead[L]{\small Gauthier Baudelet - 0472522463}
\fancyhead[R]{\small Calcul littéral avancé}
\fancyfoot[L]{\tiny Usage strictement personnel -- Diffusion interdite}
\fancyfoot[C]{\small\thepage}
\renewcommand{\headrulewidth}{0.4pt}

% ---------- Mise en forme ----------
\setlength{\parindent}{0pt}
\setlength{\parskip}{0.5em}
\setlist[itemize]{itemsep=2pt,topsep=2pt}
\setlist[enumerate]{itemsep=2pt,topsep=2pt}

\newtcolorbox{definition}{breakable,colback=blue!4,colframe=blue!60!black,
  fonttitle=\bfseries,title={Définition}}
\newtcolorbox{propriete}[1][Propriétés]{breakable,colback=green!5!white,colframe=green!45!black,
  fonttitle=\bfseries,title={#1}}
\newtcolorbox{attention}{breakable,colback=red!4,colframe=red!60!black,
  fonttitle=\bfseries,title={Attention}}
\newtcolorbox{methode}{breakable,colback=orange!6,colframe=orange!70!black,
  fonttitle=\bfseries,title={Méthode}}

\newcommand{\dec}{{,}} % virgule décimale en mode math

\title{\textbf{Calcul littéral avancé -- Produits remarquables et factorisation}\\[2pt]\large Synthèse de cours et exercices (3e secondaire)}
\author{}
\date{}

\begin{document}
\maketitle
\thispagestyle{fancy}
\vspace{-2cm}
\textbf{Propriété intellectuelle :} Ce document est rédigé et mis à disposition exclusivement pour les élèves inscrits aux cours de Gauthier Baudelet. Toute reproduction, mise en ligne (Google Drive partagé, réseaux sociaux, plateformes publiques) ou transmission à des tiers, même à titre gratuit, est strictement interdite.
% =====================================================
\section{Les produits remarquables}
% =====================================================

Un produit remarquable est une formule de calcul algébrique qui permet de développer rapidement certains produits de polynômes sans passer par la double distributivité terme à terme.

\subsection{Les trois formules fondamentales}

\begin{propriete}[Formules des produits remarquables]
Quels que soient les nombres réels $a$ et $b$ :
\begin{enumerate}[label=\textbf{(\arabic*)}]
  \item \textbf{Carré d'une somme :} 
  \[
  (a + b)^2 = a^2 + 2ab + b^2
  \]
  \item \textbf{Carré d'une différence :} 
  \[
  (a - b)^2 = a^2 - 2ab + b^2
  \]
  \item \textbf{Produit de la somme par la différence (ou différence de carrés) :} 
  \[
  (a + b)(a - b) = a^2 - b^2
  \]
\end{enumerate}
\end{propriete}

\subsection{Illustration géométrique du carré d'une somme}

Pour comprendre pourquoi $(a+b)^2 = a^2 + 2ab + b^2$, visualisons l'aire d'un grand carré de côté $(a+b)$ :

\begin{center}
\begin{tikzpicture}[scale=0.9]
  % Grand carré
  \draw[thick] (0,0) rectangle (5,5);
  \draw[dashed, thick, blue!60!black] (2,0) -- (2,5);
  \draw[dashed, thick, blue!60!black] (0,3) -- (5,3);
  
  % Aires intérieures et étiquettes
  \node at (1,4) {$a^2$};
  \node at (3.5,4) {$ab$};
  \node at (1,1.5) {$ab$};
  \node at (3.5,1.5) {$b^2$};
  
  % Dimensions extérieures
  \node[above] at (1,5) {$a$};
  \node[above] at (3.5,5) {$b$};
  \node[left] at (0,4) {$a$};
  \node[left] at (0,1.5) {$b$};
\end{tikzpicture}
\end{center}

L'aire totale est la somme des quatre aires intérieures : $a^2 + ab + ab + b^2 = a^2 + 2ab + b^2$.

\textbf{Exemples de développement :}
\begin{itemize}
  \item $(x + 3)^2 = x^2 + 2 \times x \times 3 + 3^2 = x^2 + 6x + 9$
  \item $(2x - 5)^2 = (2x)^2 - 2 \times (2x) \times 5 + 5^2 = 4x^2 - 20x + 25$
  \item $(4x - 3)(4x + 3) = (4x)^2 - 3^2 = 16x^2 - 9$
\end{itemize}

\begin{attention}
L'erreur classique à éviter absolument : 
\[
(a + b)^2 \neq a^2 + b^2 \quad \text{et} \quad (a - b)^2 \neq a^2 - b^2.
\]
Il manque toujours le double produit ($2ab$ ou $-2ab$) au milieu !
\end{attention}

% =====================================================
\section{La factorisation}
% =====================================================

Factoriser une expression algébrique, c'est la transformer d'une **somme** (ou différence) en un **produit** de facteurs. C'est l'opération inverse du développement.

\subsection{Méthode 1 : La mise en évidence}

\begin{methode}
On cherche un facteur commun (nombre, lettre ou parenthèse) présent dans tous les termes de l'expression, puis on applique la distributivité simple à l'envers : $ka + kb = k(a + b)$.
\end{methode}

\textbf{Exemples :}
\begin{itemize}
  \item $5x^2 - 10x = 5x(x - 2)$
  \item $3x(x+1) - 4(x+1) = (x+1)(3x - 4)$
\end{itemize}

\subsection{Méthode 2 : Utiliser les produits remarquables}

On lit les identités remarquables \textbf{de droite à gauche} :
\begin{enumerate}
  \item $a^2 + 2ab + b^2 = (a + b)^2$
  \item $a^2 - 2ab + b^2 = (a - b)^2$
  \item $a^2 - b^2 = (a + b)(a - b)$
\end{enumerate}

\textbf{Exemples :}
\begin{itemize}
  \item $x^2 + 8x + 16 = x^2 + 2 \times x \times 4 + 4^2 = (x + 4)^2$
  \item $9x^2 - 25 = (3x)^2 - 5^2 = (3x - 5)(3x + 5)$
\end{itemize}

% =====================================================
\section{Récapitulatif}
% =====================================================

\begin{tcolorbox}[colback=gray!6,colframe=gray!60!black,title={\bfseries À retenir}]
\textbf{Développer} $\rightarrow$ transformer un produit en somme.\\
\textbf{Factoriser} $\rightarrow$ transformer une somme en produit.

\textbf{Les 3 identités vitales :}
\[
(a+b)^2 = a^2 + 2ab + b^2 \qquad | \qquad (a-b)^2 = a^2 - 2ab + b^2 \qquad | \qquad (a+b)(a-b) = a^2 - b^2
\]
\end{tcolorbox}

\newpage
% =====================================================
\section{Exercices}
% =====================================================

\subsection{Série A -- Développer à l'aide des produits remarquables}

\begin{enumerate}[label=\textbf{\arabic*.}]

\item Développe et réduis les expressions suivantes :
\begin{enumerate}[label=\alph*)]
  \item $(x + 4)^2$ \qquad
  \item $(3x - 1)^2$ \qquad
  \item $(5 + 2x)^2$ \qquad
  \item $(x - 7)(x + 7)$
  \item $(2x - 3)(2x + 3)$ \qquad
  \item $(-x + 4)^2$ \qquad
  \item $(3x - 2)^-(x+1)$ \textit{(Attention aux signes !)}
\end{enumerate}

\end{enumerate}

\subsection{Série B -- Factoriser par mise en évidence}

\begin{enumerate}[label=\textbf{\arabic*.},start=2]

\item Factorise un maximum ces expressions :
\begin{enumerate}[label=\alph*)]
  \item $3x^2 + 12x$ \qquad
  \item $7x^3 - 14x^2 + 21x$ \qquad
  \item $a(x - 2) + b(x - 2)$ \qquad
  \item $4x(2x - 1) - 3(2x - 1)$
\end{enumerate}

\end{enumerate}

\subsection{Série C -- Factoriser à l'aide des produits remarquables}

\begin{enumerate}[label=\textbf{\arabic*.},start=3]

\item Factorise en reconnaissant un produit remarquable :
\begin{enumerate}[label=\alph*)]
  \item $x^2 + 10x + 25$ \qquad
  \item $x^2 - 9$ \qquad
  \item $4x^2 - 12x + 9$ \qquad
  \item $16x^2 - 49$ \qquad
  \item $x^2 - 2x + 1$
\end{enumerate}

\end{enumerate}

\subsection{Défi bonus}

\begin{enumerate}[label=\alph*)]
  \item Développe l'expression : $(x + 1)^3$ en utilisant $(x+1)^2(x+1)$.
  \item Factorise l'expression suivante par double étape (mise en évidence puis identité) : $2x^3 - 18x$.
\end{enumerate}

\newpage
% =====================================================
\section{Corrigés}
% =====================================================

\subsection*{Série A}
\textbf{1.} 
a) $x^2 + 8x + 16$ \quad 
b) $9x^2 - 6x + 1$ \quad 
c) $25 + 20x + 4x^2$ \quad 
d) $x^2 - 49$ \quad 
e) $4x^2 - 9$ \quad 
f) $x^2 - 8x + 16$ \quad 
g) $(9x^2 - 12x + 4) - (x + 1) = 9x^2 - 13x + 3$.

\subsection*{Série B}
\textbf{2.} 
a) $3x(x + 4)$ \quad 
b) $7x(x^2 - 2x + 3)$ \quad 
c) $(x - 2)(a + b)$ \quad 
d) $(2x - 1)(4x - 3)$.

\subsection*{Série C}
\textbf{3.} 
a) $(x + 5)^2$ \quad 
b) $(x - 3)(x + 3)$ \quad 
c) $(2x - 3)^2$ \quad 
d) $(4x - 7)(4x + 7)$ \quad 
e) $(x - 1)^2$.

\subsection*{Défi bonus}
a) $(x^2 + 2x + 1)(x+1) = x^3 + 3x^2 + 3x + 1$.\\
b) $2x(x^2 - 9) = 2x(x - 3)(x + 3)$.

\end{document}